\section{2. Gestión del proyecto con
Scrum}\label{gestiuxf3n-del-proyecto-con-scrum}

\subsection{2.1 Acuerdos de trabajo}\label{acuerdos-de-trabajo}

Los acuerdos completos están en CONTRIBUTING.md del repositorio. En
resumen:

\textbf{• Comunicación:} grupo de WhatsApp del equipo para lo
urgente y GitHub (issues y Pull Requests) para todo lo técnico. Tiempo
máximo de respuesta: 24 horas en días hábiles.

\textbf{• Reuniones:} se hacen de manera presencial o por llamada en el
grupo de WhatsApp: Daily breve de lunes a viernes y los eventos de Scrum
al inicio y al cierre de cada sprint de dos semanas. Lo que se decide en
una reunión se registra en GitHub (issue, minuta o documento).

\textbf{• Decisiones:} por consenso; si no se logra, decide el Product
Owner en temas de alcance y prioridad, y el Scrum Master en temas de
proceso.

\textbf{• Ramas:} seguimos GitHub Flow. main está protegida y cada
historia o tarea se trabaja en su propia rama, creada desde main, con la
convención \texttt{tipo/número-descripción}, donde el tipo es
\texttt{feat}, \texttt{fix}, \texttt{docs}, \texttt{test} o \texttt{ci} y
el número es el issue que atiende (por ejemplo,
\texttt{fix/9-version-stakeholders} para la observación B11). Cada rama
se integra a main mediante un Pull Request que cierra su issue y tiene al
menos una revisión de otro integrante; después se borra. Las ramas por
integrante (\texttt{rama/<NOMBRE>}) se usaron hasta la v0.2 y ya no
reciben trabajo nuevo.

\textbf{• Commits:} en español y en presente, con el formato
«Rol: mensaje» y referencia a la historia cuando aplique (por ejemplo,
«Developers: implementa recepción con checklist (HU-01)»).

\textbf{• Documentos:} el Documento de Ingeniería se escribe en LaTeX,
un archivo por sección en \texttt{docs/documento/secciones/}, para que
git pueda fusionar los cambios de varios roles. El entregable se publica en
\texttt{docs/entregas/} en PDF y en Word (.docx) con el nombre
\texttt{IS1-Fase1-EquipoB-v<versión>} (por ejemplo, IS1-Fase1-EquipoB-v0.1.pdf).

\textbf{• Incumplimiento:} se registra en la retrospectiva; a la segunda
falta en el mismo sprint, el Scrum Master lo informa al profesor.

\subsection{2.2 Definición de Listo (DoR) y Definición de Terminado
(DoD)}\label{definiciuxf3n-de-listo-dor-y-definiciuxf3n-de-terminado-dod}

{\def\LTcaptype{none} % do not increment counter
\begin{longtable}[]{@{}
  >{\raggedright\arraybackslash}p{(\linewidth - 2\tabcolsep) * \real{0.5001}}
  >{\raggedright\arraybackslash}p{(\linewidth - 2\tabcolsep) * \real{0.4999}}@{}}
\toprule\noalign{}
\begin{minipage}[b]{\linewidth}\raggedright
\textbf{Definición de Listo}
\end{minipage} & \begin{minipage}[b]{\linewidth}\raggedright
\textbf{Definición de Terminado}
\end{minipage} \\
\midrule\noalign{}
\endhead
\bottomrule\noalign{}
\endlastfoot
• La historia sigue el formato Como\ldots{} quiero\ldots{} para\ldots{}
y tiene ID e issue en el tablero.

• Tiene al menos 2 criterios de aceptación en formato
Dado/Cuando/Entonces.

• El equipo la estimó en puntos y cabe en un sprint (8 puntos o menos).

• Sus reglas de negocio (RN) y requisitos (RF/RNF) relacionados están
identificados.

• Sus dependencias (otras historias, datos, wireframe) están resueltas o
planificadas.

• El Product Owner la priorizó (MoSCoW). & • El código está en main
mediante un Pull Request revisado por otro integrante.

• Cumple todos sus criterios de aceptación y el Product Owner lo
verificó.

• Tiene pruebas automatizadas y la integración continua pasa sin
errores.

• Se probó solo con datos ficticios; no hay datos personales reales ni
contraseñas en el repositorio.

• La documentación afectada (SRS, modelos, README) está actualizada.

• Está desplegada en el ambiente de pruebas (desde el Sprint 4). \\
\end{longtable}
}

\subsection{2.3 Resumen de sprints}\label{resumen-de-sprints}

{\def\LTcaptype{none} % do not increment counter
\begin{longtable}[]{@{}
  >{\raggedright\arraybackslash}p{(\linewidth - 10\tabcolsep) * \real{0.0856}}
  >{\raggedright\arraybackslash}p{(\linewidth - 10\tabcolsep) * \real{0.1495}}
  >{\raggedright\arraybackslash}p{(\linewidth - 10\tabcolsep) * \real{0.3270}}
  >{\raggedright\arraybackslash}p{(\linewidth - 10\tabcolsep) * \real{0.1601}}
  >{\raggedright\arraybackslash}p{(\linewidth - 10\tabcolsep) * \real{0.1389}}
  >{\raggedright\arraybackslash}p{(\linewidth - 10\tabcolsep) * \real{0.1389}}@{}}
\toprule\noalign{}
\begin{minipage}[b]{\linewidth}\raggedright
\textbf{Sprint}
\end{minipage} & \begin{minipage}[b]{\linewidth}\raggedright
\textbf{Fechas}
\end{minipage} & \begin{minipage}[b]{\linewidth}\raggedright
\textbf{Objetivo del sprint}
\end{minipage} & \begin{minipage}[b]{\linewidth}\raggedright
\textbf{¿Se cumplió? (sí / parcial / no)}
\end{minipage} & \begin{minipage}[b]{\linewidth}\raggedright
\textbf{Puntos comprometidos}
\end{minipage} & \begin{minipage}[b]{\linewidth}\raggedright
\textbf{Puntos terminados}
\end{minipage} \\
\midrule\noalign{}
\endhead
\bottomrule\noalign{}
\endlastfoot
S0 & \pendiente{} & Conocer al cliente, levantar los requisitos
iniciales y preparar el repositorio y el documento base. &
Sí & --- & --- \\
S1 & \pendiente{} & SRS en borrador (4.1--4.7) para validarlo con el
cliente. & Sí: el borrador se valida con el cliente en la Review 1 &
0 & 0 \\
\end{longtable}
}

Los sprints 0 y 1 fueron de especificación: su objetivo no incluyó
historias de usuario, así que no hay puntos comprometidos ni terminados
y las 30 historias siguen en estado «Pendiente» (4.5). La velocidad se
calcula a partir del Sprint 2, con los puntos de las historias
terminadas según la DoD.

\subsubsection{Evidencia de los eventos}\label{evidencia-de-los-eventos}

{\def\LTcaptype{none} % do not increment counter
\begin{longtable}[]{@{}
  >{\raggedright\arraybackslash}p{(\linewidth - 8\tabcolsep) * \real{0.0962}}
  >{\raggedright\arraybackslash}p{(\linewidth - 8\tabcolsep) * \real{0.2136}}
  >{\raggedright\arraybackslash}p{(\linewidth - 8\tabcolsep) * \real{0.2137}}
  >{\raggedright\arraybackslash}p{(\linewidth - 8\tabcolsep) * \real{0.3162}}
  >{\raggedright\arraybackslash}p{(\linewidth - 8\tabcolsep) * \real{0.1604}}@{}}
\toprule\noalign{}
\begin{minipage}[b]{\linewidth}\raggedright
\textbf{Sprint}
\end{minipage} & \begin{minipage}[b]{\linewidth}\raggedright
\textbf{Minuta de Review}
\end{minipage} & \begin{minipage}[b]{\linewidth}\raggedright
\textbf{Retrospectiva}
\end{minipage} & \begin{minipage}[b]{\linewidth}\raggedright
\textbf{Acciones de mejora acordadas}
\end{minipage} & \begin{minipage}[b]{\linewidth}\raggedright
\textbf{¿Se aplicaron?}
\end{minipage} \\
\midrule\noalign{}
\endhead
\bottomrule\noalign{}
\endlastfoot
S0 & \pendiente{} & \pendiente{} & --- & --- \\
S1 & \ruta{docs/minutas/2026-11-10-review-1.md} & \pendiente{} & --- & --- \\
\end{longtable}
}

\subsection{2.4 Registro de riesgos}\label{registro-de-riesgos}

{\def\LTcaptype{none} % do not increment counter
\begin{longtable}[]{@{}
  >{\raggedright\arraybackslash}p{(\linewidth - 10\tabcolsep) * \real{0.0856}}
  >{\raggedright\arraybackslash}p{(\linewidth - 10\tabcolsep) * \real{0.2990}}
  >{\raggedright\arraybackslash}p{(\linewidth - 10\tabcolsep) * \real{0.1282}}
  >{\raggedright\arraybackslash}p{(\linewidth - 10\tabcolsep) * \real{0.1069}}
  >{\raggedright\arraybackslash}p{(\linewidth - 10\tabcolsep) * \real{0.2735}}
  >{\raggedright\arraybackslash}p{(\linewidth - 10\tabcolsep) * \real{0.1067}}@{}}
\toprule\noalign{}
\begin{minipage}[b]{\linewidth}\raggedright
\textbf{ID}
\end{minipage} & \begin{minipage}[b]{\linewidth}\raggedright
\textbf{Riesgo}
\end{minipage} & \begin{minipage}[b]{\linewidth}\raggedright
\textbf{Probabilidad}
\end{minipage} & \begin{minipage}[b]{\linewidth}\raggedright
\textbf{Impacto}
\end{minipage} & \begin{minipage}[b]{\linewidth}\raggedright
\textbf{Respuesta y responsable}
\end{minipage} & \begin{minipage}[b]{\linewidth}\raggedright
\textbf{Estado}
\end{minipage} \\
\midrule\noalign{}
\endhead
\bottomrule\noalign{}
\endlastfoot
R-01 & Un integrante abandona el curso a mitad del semestre & Media &
Alto & Mitigar: conocimiento compartido mediante revisión cruzada de PR;
responsable: SM & Abierto \\
R-02 & La operación sin conexión y la sincronización posterior (RNF-08,
RNF-09) resultan más complejas de lo estimado & Alta & Alto & Mitigar:
prototipo técnico en el Sprint 1 y limitar lo offline a recepción y
diagnóstico; responsable: equipo de desarrollo & Abierto \\
R-03 & El cliente no está disponible para validar requisitos o asistir a
las Reviews & Media & Alto & Mitigar: agendar sesiones con anticipación
y validar por correo cuando no pueda asistir; responsable: PO &
Abierto \\
R-04 & Cambios de alcance frecuentes por parte del cliente & Media &
Medio & Mitigar: registrar cada solicitud de cambio (2.5) con análisis
de impacto; responsable: PO & Abierto \\
R-05 & El servicio gratuito de correo limita el envío de enlaces y
códigos 2FA & Media & Alto & Mitigar: probar los límites del plan
gratuito en el Sprint 1 y dejar SMS como opción (RF-02.3); responsable:
equipo de desarrollo & Abierto \\
R-06 & Las fotografías de recepción saturan el almacenamiento del
hospedaje gratuito & Media & Medio & Mitigar: comprimir las imágenes en
el dispositivo y fijar un tamaño máximo por foto; responsable: equipo de
desarrollo & Abierto \\
R-07 & Poca experiencia del equipo con FastAPI y PWA & Alta & Medio &
Mitigar: capacitación en S0--S1 y trabajo en pares en las primeras
historias; responsable: SM & Abierto \\
R-08 & Conflictos al editar en paralelo el documento .docx en ramas
distintas & Alta & Medio & Evitar: el documento se migró a LaTeX con un
archivo por sección; responsable: SM & Cerrado \\
R-09 & Se publican datos personales reales en el repositorio público &
Media & Alto & Evitar: solo datos ficticios y revisión de datos
personales en cada PR; .gitignore excluye .env y credenciales;
responsable: todos & Abierto \\
R-10 & Código o texto generado con IA que el equipo no entiende o no
revisó & Media & Alto & Mitigar: declarar el uso de IA en el documento, revisión de
cada PR por otro integrante y pruebas automatizadas como evidencia;
responsable: SM & Abierto \\
R-11 & Diferencias entre SQLite (pruebas) y PostgreSQL (staging) ocultan
defectos & Media & Medio & Mitigar: cargar el esquema en PostgreSQL
en la integración continua desde que exista el backend (S2);
responsable: QA Lead & Abierto \\
R-12 & Sin ambiente de pruebas no se puede mostrar el flujo núcleo al
cliente & Media & Alto & Mitigar: decidir el hospedaje del ambiente de pruebas antes del Sprint 4;
responsable: Developers & Abierto \\
\end{longtable}
}

\subsection{2.6 Declaración de uso de inteligencia
artificial}\label{declaracion-de-uso-de-ia}

El equipo usó un asistente de inteligencia artificial para pasar el
documento a LaTeX, atender las observaciones de la Review 1 y organizar el
repositorio. La tabla indica en qué se usó, cómo se verificó el resultado
y quién responde por él.

{\small
\begin{longtable}[]{@{}
  >{\raggedright\arraybackslash}p{0.13\linewidth}
  >{\raggedright\arraybackslash}p{0.21\linewidth}
  >{\raggedright\arraybackslash}p{0.20\linewidth}
  >{\raggedright\arraybackslash}p{0.22\linewidth}
  >{\raggedright\arraybackslash}p{0.13\linewidth}@{}}
\toprule
\textbf{Herramienta} & \textbf{Sección o artefacto} & \textbf{Para qué se usó} & \textbf{Cómo se verificó} & \textbf{Responsable} \\
\midrule
\endhead
\bottomrule
\endlastfoot
Claude Code (Anthropic) & Documento en LaTeX (\ruta{docs/documento/}):
migración desde el Word v0.1, plantilla, preámbulo y
\ruta{generar-docx.py} & Pasar a LaTeX el contenido que el equipo ya
había redactado en Word, con un archivo por sección & El documento
compila con pdflatex sin errores & Abdías Daniel Chávez Cruz \\
Claude Code (Anthropic) & Correcciones B11 y B5: 3.4, Figura 3.1,
portada, control de versiones, CONTRIBUTING.md, README, 2.1 y Anexo F &
Redactar los cambios que pide cada observación & Compilación sin errores
y comprobación de que no quedan los stakeholders retirados ni las ramas
por rol & Abdías Daniel Chávez Cruz \\
Claude Code (Anthropic) & Corrección B3: Anexo A (guion y minuta de la
entrevista) y fuentes de las reglas de 4.7 & Redactar el guion y la
minuta, y ligar cada regla con su hallazgo & Cada regla de entrevista de
4.7 cita un hallazgo de la minuta. La minuta está pendiente de validar
por el Product Owner contra sus notas de la entrevista & Abdías Daniel
Chávez Cruz \\
Claude Code (Anthropic) & Corrección B4: issues de las 30 historias,
números en 4.5, 2.3, minuta de la Review 1 y 12.2 & Registrar los issues
a partir de 4.5, resumir la retroalimentación del profesor y contar
commits, PR y revisiones & Datos comprobados contra GitHub y contra
\ruta{informacion/correcionesContexto.md} & Abdías Daniel Chávez Cruz \\
Claude Code (Anthropic) & Propuesta de alcance del Release v0.1.0
(observación B2) & Redactar la propuesta a partir de 4.5 y del
anteproyecto & Sumas de historias y puntos comprobadas contra 4.5.
Pendiente de revisión por el Product Owner y de negociación con el
cliente & Abdías Daniel Chávez Cruz \\
Claude Code (Anthropic) & Versiones v0.2.1 a v0.2.3, etiquetas de Git,
issues, ramas, Pull Requests, carpeta \ruta{informacion/} y esta
declaración & Aplicar los acuerdos de CONTRIBUTING.md en el repositorio &
Estado del repositorio comprobado en GitHub después de cada cambio &
Abdías Daniel Chávez Cruz \\
\pendiente{} & Contenido original de las secciones 1 a 4, redactado en
Word en octubre de 2026 & \pendiente{} & \pendiente{} & Cada autor declara
su parte \\
\end{longtable}
}

\textbf{Datos personales.} No se ingresaron datos de clientes en la
herramienta: el proyecto solo usa datos ficticios. El asistente trabajó
sobre el contenido del repositorio, que incluye los nombres y las
matrículas de los integrantes (sección 1.1).

\textbf{Concentración del trabajo.} Todo el uso del asistente estuvo a
cargo de un solo integrante. Para repartir ese conocimiento, el Tech Lead
explicará al equipo los cambios hechos en una sesión interna con minuta
(observación B6, fecha límite 18/11/2026).
